\section{From Pairwise Metrics to Multi-Task Mergeability}
\label{sec:multitask_prediction}

A natural question arises: can pairwise mergeability metrics predict the success of merging more than two models? We investigate whether the linear predictors learned from pairwise data generalize to multi-task merging scenarios.

\subsection{Experimental Setup}

\paragraph{Subset Sampling}
We sample 200 random subsets, each containing 5 tasks from our benchmark of 20 tasks. This yields subsets with $\binom{5}{2} = 10$ pairwise interactions each, providing sufficient diversity while remaining computationally tractable.

\paragraph{Predicted Mergeability}
For each subset $S = \{t_1, t_2, t_3, t_4, t_5\}$, we define the \emph{predicted subset mergeability} as the average of pairwise predicted mergeabilities:
\begin{equation}
    \hat{M}(S) = \frac{1}{10} \sum_{i < j} \sum_{m=1}^{28} w_m \cdot x_m(t_i, t_j)
\end{equation}
where $w_m$ are the learned L1 LOTO coefficients for a given merging method, and $x_m(t_i, t_j)$ is the value of metric $m$ for the pair $(t_i, t_j)$. This formulation assumes that multi-task mergeability can be approximated by aggregating pairwise interactions.

\paragraph{Actual Performance}
We perform the actual 5-way merge for all 200 subsets using each merging method (TSV, TIES, Weight Averaging, Task Arithmetic, DARE). We measure the normalized accuracy averaged across the 5 constituent tasks.

\paragraph{Evaluation Protocol}
We rank the 200 subsets by predicted mergeability and compare the top 10 (highest predicted) against the bottom 10 (lowest predicted). If pairwise metrics generalize to multi-task settings, the top-predicted subsets should achieve higher actual performance than the bottom-predicted subsets.

\subsection{Results}

Table~\ref{tab:multitask_prediction} presents the comparison between predicted and actual performance for the extreme subsets.

\begin{table}[h]
\centering
\caption{Top 10 vs.\ Bottom 10 predicted subsets: comparison of actual merging performance. All five methods correctly identify that top-predicted subsets outperform bottom-predicted subsets, with TSV showing the largest and most reliable gap.}
\label{tab:multitask_prediction}
\begin{tabular}{@{}lccccr@{}}
\toprule
\textbf{Method} & \textbf{Top 10 Actual} & \textbf{Bottom 10 Actual} & \textbf{$\Delta$ Accuracy} & \textbf{Correct?} \\
\midrule
TSV & 0.958 & 0.923 & \textbf{+3.5\%} & \checkmark \\
Task Arithmetic & 0.888 & 0.859 & +2.9\% & \checkmark \\
DARE & 0.856 & 0.838 & +1.7\% & \checkmark \\
TIES & 0.838 & 0.821 & +1.7\% & \checkmark \\
Weight Averaging & 0.857 & 0.844 & +1.2\% & \checkmark \\
\bottomrule
\end{tabular}
\end{table}

\subsection{Discussion}

\paragraph{Pairwise Metrics Generalize to Multi-Task Merging}
All five merging methods correctly distinguish high-mergeability subsets from low-mergeability subsets based solely on pairwise metric predictions. This validates our core hypothesis: pairwise interactions capture meaningful information about multi-task compatibility.

\paragraph{TSV Shows Strongest Signal}
TSV exhibits the largest performance gap (+3.5\%) between top and bottom predicted subsets. This suggests that TSV's merging mechanism---which operates on singular value decompositions---is particularly well-captured by our pairwise metrics, many of which measure subspace and gradient similarities. The linear superposition assumption underlying our prediction (averaging pairwise scores) aligns well with TSV's mathematical formulation.

\paragraph{Practical Implications}
These results have practical implications for model merging at scale. When faced with a large pool of candidate models, practitioners can use pairwise metrics to efficiently screen subsets before committing to expensive multi-way merges. Even for methods where the correlation is weaker (Weight Averaging, TIES), the prediction reliably identifies subsets to avoid.

\paragraph{Limitations}
While the direction of prediction is correct for all methods, the magnitude of the gap varies substantially. For Weight Averaging and TIES, the +1--2\% improvement may not justify the computational cost of pairwise metric computation in all scenarios. Additionally, our aggregation strategy (simple averaging) may not capture higher-order interactions that emerge when merging more than two models.
